\bookchapter{Chapitre 17 : Interaction entre le smartphone et d'autres appareils}

\chapterintro{Le smartphone que vous tenez en main est en réalité une super-télécommande et un centre de données. Il ne sert pas qu'à passer des appels, il peut aussi connecter ordinateur, TV, montre, enceinte, voiture, pour que ces appareils ne soient plus des îles isolées. Ce chapitre parle de comment les faire « faire équipe ».}

\booksection{I. Collaboration avec l'ordinateur : le smartphone et l'ordinateur ne sont plus deux mondes}

Honnêtement, ce que je détestais avant, c'était de prendre une photo sur le smartphone, vouloir la retoucher sur l'ordinateur, et devoir trouver un câble de données, attendre le transfert. J'ai découvert ensuite qu'en réalité, aucun câble n'était nécessaire, une bonne intégration est dix fois plus rapide qu'un câble.

\begin{itemize}
  \item \textbf{Écosystème Apple (AirDrop + presse-papiers universel)}: si vous utilisez un iPhone avec un Mac, AirDrop est la première fonction à apprendre. Sur le smartphone, appuyez sur Partager AirDrop sélectionnez votre Mac, la photo passe instantanément, sans données, sans WiFi, utilise Bluetooth et WiFi Direct. Il y a aussi une fonction facilement négligée appelée « presse-papiers universel » — copiez un texte sur le smartphone, collez-le directement sur le Mac, c'est comme de la magie.
  \item \textbf{Android + Windows (Link to Windows)}: Microsoft a une application appelée « Link to Windows », intégrée aux téléphones Samsung, les autres téléphones Android peuvent la télécharger depuis la boutique d'applications. Après jumelage, vous pouvez directement voir les notifications du téléphone sur l'ordinateur, répondre aux SMS, parcourir les photos, même passer et recevoir des appels. Personnellement, je l'utilise depuis deux ans, en réunion pas besoin de regarder constamment le téléphone, une notification qui apparaît sur l'écran de l'ordinateur suffit.
  \item \textbf{Collaboration multi-écran de marque}: le « Multi-Screen Collaboration » de Huawei et le « Xiaomi MiaXiang » de Xiaomi vont plus loin — ils projettent directement l'écran du smartphone sur l'ordinateur, vous pouvez utiliser la souris pour faire glisser des fichiers du smartphone à l'ordinateur, et aussi surfer sur TikTok, répondre à WeChat sur l'ordinateur. Copier-coller entre appareils, glisser-déposer de fichiers, l'expérience est très fluide.
\end{itemize}

\booksection{II. Casting vers la TV : tout le monde regarde ensemble}

Pour le Nouvel An, ma cousine voulait diffuser sur la TV la vidéo de son voyage de fin d'études qui était sur son téléphone. Après avoir bidouillé pendant longtemps, elle a découvert que c'était en réalité très simple :

\begin{itemize}
  \item \textbf{iPhone TV}: faites glisser vers le bas le centre de contrôle depuis le coin supérieur droit, appuyez sur l'icône avec deux cadres superposés (Screen Mirroring), trouvez le nom de votre TV, un appui et c'est connecté. Condition préalable : la TV et le téléphone doivent être sur le même WiFi, et la TV doit prendre en charge AirPlay (la plupart des TV intelligentes le prennent en charge).
  \item \textbf{Android TV}: les appellations varient, Huawei l'appelle « projection sans fil », Xiaomi « projection », OPPO/vivo « interaction multi-écran ». Généralement dans le panneau des raccourcis, sinon cherchez « projection » dans les réglages.
  \item \textbf{Application vidéo avec casting direct}: si vous voulez juste regarder iQiyi, Bilibili, YouTube sur la TV, c'est encore plus simple — dans l'interface de lecture, cherchez une petite icône rectangulaire avec un symbole WiFi, appuyez dessus, sélectionnez la TV, la vidéo passe sur la TV. Le téléphone peut faire autre chose, la vidéo continue sur la TV, sans interférence.
  \item \textbf{Que faire si le casting échoue}: 90 \% des cas, c'est parce que vous n'êtes pas sur le même WiFi. Les 10 \% restants, c'est que la fonction de casting de la TV n'est pas activée, cherchez l'interrupteur « interaction multi-écran » ou « AirPlay » dans les réglages de la TV.
\end{itemize}

\booksection{III. Avec la montre/bracelet connecté : un coach santé au poignet}

J'ai offert un bracelet à ma mère, au début elle disait « je ne fais pas de sport, à quoi ça sert ? ». Après deux semaines d'utilisation, elle a découvert qu'elle ne dormait que cinq heures par jour, et a commencé à ajuster activement son rythme de vie. C'est la valeur la plus pratique des montres/bracelets — transformer les données de santé invisibles en chiffres visibles.

\begin{itemize}
  \item \textbf{Filtre de notifications}: quand le téléphone vibre sans arrêt dans le sac ou la poche, levez le poignet pour voir — c'est un WeChat du patron ou une promo Pinduoduo, jugez en une seconde, les messages qui ne méritent pas de sortir le téléphone, laissez-les passer.
  \item \textbf{Synchronisation des données de santé}: la montre surveille en continu 24h/24 le rythme cardiaque, l'oxygène sanguin, le sommeil, le nombre de pas, ces données sont automatiquement agrégées dans l'application Santé sur le smartphone. Avec le temps, vous voyez vos tendances — ce mois-ci ai-je trop peu marché, cette semaine la qualité du sommeil s'est-elle améliorée.
  \item \textbf{Synchronisation carte de transport/carte d'accès}: les cartes de bus et cartes d'accès configurées sur le smartphone sont automatiquement synchronisées sur la montre chez la plupart des marques. Dans le métro, pas besoin de sortir le téléphone, le poignet s'approche du portail et c'est passé, on a vraiment l'impression que « le futur est arrivé ».
  \item \textbf{Trouver mon téléphone, l'outil magique}: la montre a généralement une fonction « trouver mon téléphone », un appui et le téléphone sonne. Personnellement, j'utilise cette fonction bien plus que prévu — je pose le téléphone n'importe où et oublie où il est.
\end{itemize}

\booksection{IV. Avec la maison connectée : le smartphone est le « commutateur général » de la maison}

La maison connectée paraît mystérieuse, mais en réalité, ça se résume à trois choses : contrôle à distance par smartphone, définition de règles d'automatisation, contrôle vocal.

Laissez-moi vous parler de mes propres usages pratiques, vous comprendrez à quel point c'est utile :

\begin{itemize}
  \item \textbf{Climatisation}: en été, sur le chemin du retour, j'allume la clim à distance avec le smartphone, en entrant c'est déjà frais. Avec une prise adaptateur pour clim à une centaine de yuans, c'est possible, pas besoin de changer de clim.
  \item \textbf{Lumière}: avant de dormir, au lit, trop paresseux pour se lever, appuyez sur le smartphone ou dites « Hey Siri/Xiao Ai, éteins la lumière ».
  \item \textbf{Robot aspirateur}: avant de sortir, appuyez sur « commencer le nettoyage » sur le smartphone, en rentrant le sol est propre. Peut aussi être programmé pour aspirer tous les jours à heure fixe.
  \item \textbf{Caméra}: avec des personnes âgées, des enfants ou des animaux à la maison, vous pouvez à tout moment jeter un œil à la situation depuis le smartphone.
\end{itemize}

Quel écosystème choisir ? En résumé :

\begin{itemize}
  \item Si vous utilisez un iPhone privilégiez les appareils prenant en charge Apple HomeKit (un peu plus chers mais meilleure protection de la vie privée).
  \item Beaucoup d'appareils et rapport qualité-prix l'écosystème Xiaomi Mi Home est le plus complet, lumières, prises, aspirateurs, il y a de tout.
  \item Si vous utilisez un Huawei privilégiez HarmonyOS Link, meilleure expérience avec les appareils Huawei.
\end{itemize}

\booksection{V. Écouteurs Bluetooth et enceintes : la liberté sans fil}

Passer des écouteurs filaires aux écouteurs Bluetooth a été pour moi l'amélioration la plus notable de l'expérience smartphone. Plus jamais un sac rempli de fils emmêlés, plus jamais de fil accroché à une poignée de porte.

\begin{itemize}
  \item \textbf{Jumelage simple}: les écouteurs entrent généralement automatiquement en mode de jumelage à la première mise sous tension (clignotement), le smartphone avec Bluetooth activé les détecte, un appui et c'est connecté. Ensuite, à chaque fois que vous les sortez, connexion automatique, rien à faire.
  \item \textbf{Avantages de la même marque}: les AirPods avec iPhone affichent automatiquement une fenêtre pop-up avec le niveau de batterie, et peuvent basculer de manière transparente entre plusieurs appareils Apple. Les FreeBuds Huawei avec téléphone Huawei, les écouteurs Xiaomi avec téléphone Xiaomi ont des fenêtres pop-up et optimisations similaires.
  \item \textbf{Enceinte Bluetooth}: lors de réunions familiales, de camping en plein air, une enceinte Bluetooth peut doubler l'ambiance. Une fois connectée, le smartphone devient votre table DJ.
  \item \textbf{Attention à la latence}: en jouant avec des écouteurs Bluetooth, il peut y avoir une légère latence (le coup de feu est entendu une demi-seconde plus tard). Si vous êtes sensible à la latence, utilisez des écouteurs filaires ou achetez des écouteurs Bluetooth prenant en charge le « mode faible latence ».
\end{itemize}

\begin{tcbinfo}{Technologie UWB (ultra large bande) — un « sens de l'espace » plus précis que le Bluetooth}
Outre le Bluetooth, les nouveaux smartphones cachent également une puce UWB (Ultra-Wideband, ultra large bande). Sa capacité est\textbf{la mesure de distance précise et le positionnement spatial}— elle peut calculer à quelle distance se trouve un autre appareil et dans quelle direction, avec une précision de l'ordre du centimètre. Le Bluetooth peut seulement vous dire « à proximité », l'UWB peut vous dire « à 1,2 m en avant à gauche ». Applications courantes :

\begin{itemize}
  \item \textbf{Précision AirTag/Localiser}: Lorsque l'iPhone localise un AirTag, l'affichage « 7 m, tournez à gauche » avec navigation par flèche repose sur l'UWB. Le Bluetooth ordinaire ne peut indiquer que « dans ou hors de portée ».
  \item \textbf{Clé de voiture numérique}: Les iPhone et smartphones Android phares équipés de l'UWB peuvent servir de clé de voiture, déverrouillant automatiquement à l'approche et verrouillant à l'éloignement, tout en distinguant la position du conducteur de celle du passager, pour prévenir les attaques par relais.
  \item \textbf{Navigation en intérieur}: Dans les centres commerciaux et les aéroports, des balises UWB permettent une navigation « au rayon près », bien plus précise que le GPS.
  \item \textbf{Modèles pris en charge}: L'iPhone 11 et les modèles ultérieurs sont tous équipés d'une puce UWB (U1/U2) ; côté Android, certains modèles phares de Xiaomi, Samsung, vivo et OPPO en sont également équipés.
\end{itemize}

\end{tcbinfo}

\begin{tcbinfo}{LE Audio et Auracast — la nouvelle génération de normes audio Bluetooth}
L'audio Bluetooth classique (A2DP) utilisé depuis vingt ans est en train d'être remplacé par\textbf{LE Audio}(audio à faible consommation), la nouvelle norme audio de Bluetooth 5.2 et au-delà :

\begin{itemize}
  \item \textbf{Plus économe en énergie, meilleure qualité audio}: Utilise le tout nouveau codec LC3, offrant une meilleure qualité audio et une consommation plus faible que le SBC de l'ancienne norme, l'autonomie des écouteurs est nettement prolongée.
  \item \textbf{Audio double voie}: Transmission indépendante des écouteurs gauche et droit, basculement en mode mono sans déconnexion.
  \item \textbf{Audio diffusé Auracast}: La fonction phare de LE Audio. Un smartphone ou une télévision peut diffuser simultanément vers\textbf{un nombre illimité de}paires d'écouteurs — dans les salles d'embarquement, les musées, les salles de sport et les salles de classe, une seule source permet à plusieurs personnes d'écouter le même contenu avec leurs propres écouteurs.
  \item \textbf{État actuel}: En 2024-2026, les écouteurs et smartphones phares nouvellement commercialisés prennent généralement en charge LE Audio, les anciens appareils nécessitent une mise à niveau matérielle et ne peuvent pas être débloqués par logiciel. Lors de l'achat de nouveaux écouteurs, vérifiez la présence du logo LE Audio/Auracast.
\end{itemize}

\end{tcbinfo}

\booksection{VI. Avec le système embarqué : un partenaire de sécurité au volant}

Regarder son smartphone en conduisant est illégal et mortel. Mais parfois, on a vraiment besoin de navigation, de prendre un appel, de mettre de la musique. C'est là que CarPlay (Apple) et Android Auto (Android) entrent en jeu.

\begin{itemize}
  \item \textbf{Comment se connecter}: si votre voiture le prend en charge (la plupart des voitures récentes le font), branchez le smartphone avec un câble de données sur le port USB de la voiture, l'écran embarqué bascule automatiquement sur l'interface CarPlay ou Android Auto. Certaines voitures récentes prennent en charge la connexion sans fil, encore plus pratique.
  \item \textbf{Que peut-on faire}: navigation (bien meilleure que la navigation embarquée), appels téléphoniques (décroché au volant), écouter de la musique (les applications de musique du smartphone directement projetées sur la voiture), envoyer des WeChat vocaux (Siri ou Xiao Ai lit les messages et répond).
  \item \textbf{Les voitures anciennes aussi}: si votre voiture est trop ancienne et ne le prend pas en charge, achetez pour quelques centaines de yuans un écran CarPlay externe, collez-le sur le tableau de bord, l'expérience est similaire.
\end{itemize}

\booksection{VII. Sécurité et vie privée dans les interactions}

Plus les appareils sont connectés, plus la surface d'exposition aux fuites de vie privée est grande. Quelques points à noter :

\begin{itemize}
  \item Avant de projeter sur un écran en public, vérifiez d'abord si la barre de notifications du smartphone ne contient pas de messages privés, ne laissez pas toute la famille voir ce qu'elle ne devrait pas voir.
  \item Avant de donner ou vendre un appareil à quelqu'un, dissociez-le d'abord — allez dans les réglages du smartphone, trouvez l'appareil, sélectionnez « supprimer » ou « annuler le jumelage ».
  \item Pour les caméras de maison connectée, achetez des marques, avec chiffrement, n'achetez pas des caméras sans marque à quelques dizaines de yuans. Trop de cas de caméras piratées diffusées en direct par des hackers.
\end{itemize}

\begin{tcbpractice}{Essayez : projetez l'écran du smartphone sur la TV et créez une scène intelligente}
\begin{enumerate}
  \item Assurez-vous que le smartphone et la TV sont sur le même WiFi.
  \item Trouvez la fonction « Screen Mirroring » ou « AirPlay » sur la TV et activez-la.
  \item Sur le smartphone : iPhone appuyez sur « Screen Mirroring » depuis le centre de contrôle ; Android faites glisser vers le bas pour trouver « projection sans fil ».
  \item Projetez une de vos photos préférées sur la TV pour que la famille la voie.
  \item Si vous avez des appareils de maison connectée, essayez de créer un « mode départ » dans l'application smartphone : éteindre les lumières, couper la clim, lancer le robot aspirateur en un seul geste.
\end{enumerate}

\end{tcbpractice}

\begin{tcbinfo}{Cas 1 : Optimisation pratique de l'autonomie de la batterie}
\begin{itemize}
  \item \textbf{Scénario}: L'utilisateur d'un Xiaomi 14 a une mauvaise autonomie, la batterie est vide en une demi-journée.
  \item \textbf{Étapes}: (1) Réglages Batterie Consulter le classement de consommation (2) Identifier les applications en arrière-plan anormales (3) Désactiver la 5G et passer en 4G (4) Limiter l'arrière-plan (5) Activer le mode économie d'énergie (6) Retirer la coque pendant la charge pour le refroidissement.
  \item \textbf{Résultat attendu}: Autonomie prolongée de 30-50 \%, données de santé de la batterie lisibles.
\end{itemize}

\end{tcbinfo}
